\begin{abstract}
Mathematical reasoning remains a formidable obstacle for language models, given its inherent complexity and structured demands. Here, we present DeepSeekMath~7B, an advancement over DeepSeek-Coder-Base-v1.5 7B, further trained with 120B mathematics-focused tokens extracted from Common Crawl alongside text and code corpora. DeepSeekMath~7B attains a notable 51.7\% on the high-difficulty MATH benchmark, doing so independently of external solver integrations or ensemble voting strategies, and thus narrows the gap with state-of-the-art systems such as Gemini-Ultra and GPT-4. Furthermore, when employing self-consistency with 64 generated solutions, DeepSeekMath~7B achieves 60.9\% accuracy on the same benchmark. The robust mathematical reasoning skills of DeepSeekMath~arise from two pivotal advancements: Firstly, we capitalize on the vast availability of mathematical content online using a rigorously designed data filtration and selection method. Secondly, we propose Group Relative Policy Optimization (GRPO)—a novel extension of Proximal Policy Optimization (PPO)—which not only sharpens mathematical reasoning but also yields improvements in PPO's memory efficiency.
\end{abstract}

\section{Introduction}

Large language models (LLMs) have transformed the landscape of mathematical reasoning within artificial intelligence, leading to notable progress in both the quantitative reasoning benchmark \citep{MATH} and the geometry reasoning benchmark \citep{trinh2024solving}. These systems also serve as valuable collaborators for humans addressing sophisticated mathematical challenges \citep{tao}. Nonetheless, leading-edge models such as GPT-4 \citep{gpt4} and Gemini-Ultra \citep{gemini} remain unavailable to the public, while current open-source alternatives still exhibit a marked performance deficit by comparison.

This paper presents DeepSeekMath, a language model specialized in mathematical domains that demonstrates substantial advancements over existing open-source alternatives and exhibits performance on par with GPT-4 when assessed on scholarly benchmarks. Central to this achievement is the development of the DeepSeekMath Corpus, an extensive and high-quality pre-training dataset comprising 120B mathematical tokens. The corpus is constructed by leveraging a fastText-based classifier \citep{joulin2016fasttext} to extract relevant data from Common Crawl (CC). Initially, the classifier is trained with positive samples drawn from OpenWebMath \citep{openwebmath} alongside a broad selection of unrelated web pages as negative samples. With this baseline, the classifier is employed to identify additional relevant data within CC, and these results are further curated through manual annotation. The newly annotated data is then used to refine and retrain the classifier, further enhancing its effectiveness. Experimental results affirm the superior quality of the resulting corpus, as demonstrated by DeepSeekMath-Base 7B attaining a score of 64.2\% on GSM8K \citep{gsm8k} and 36.2\% on the MATH dataset \citep{MATH}, exceeding the performance of Minerva 540B \citep{minerva}. Notably, due to the multilingual nature of the DeepSeekMath Corpus, we also observe improvements on Chinese mathematical benchmarks \citep{wei2023cmath,agieval}. We view our methodology for curating mathematical data as a valuable initial contribution for further exploration and improvement within the research community.

DeepSeekMath-Base is built upon the DeepSeek-Coder-Base-v1.5 7B model \citep{deepseek-coder}, as empirical evidence suggests that initializing with a code-focused model yields superior results relative to a general large language model. Additionally, we find that mathematical pre-training confers broader benefits, enhancing performance not only on mathematics benchmarks but also improving results on the MMLU \citep{mmlu} and BBH \citep{bbh} general reasoning tasks.

Following pre-training, we conduct mathematical instruction tuning of DeepSeekMath-Base utilizing datasets incorporating chain-of-thought \citep{cot}, program-of-thought \citep{pot,pal}, and tool-augmented reasoning \citep{tora} approaches. This yields DeepSeekMath-Instruct 7B, which outperforms all other 7B-sized models and rivals the abilities of 70B parameter open-source models that have undergone instruction tuning.

Moreover, we present Group Relative Policy Optimization (GRPO), a novel reinforcement learning variant derived from Proximal Policy Optimization (PPO) \citep{schulman2017proximal}. GRPO eliminates the need for a critic network by leveraging group-based baseline estimation, leading to significant reductions in training overhead. Using only a portion of English-language instruction tuning data, GRPO significantly enhances DeepSeekMath-Instruct’s already strong performance, both in-domain (GSM8K: 82.9\% $\rightarrow$ 88.2\%, MATH: 46.8\% $\rightarrow$ 51.7\%) and on out-of-domain mathematics tasks (e.g., CMATH: 84.6\% $\rightarrow$ 88.8\%) during the reinforcement learning process. Furthermore, we propose a unifying conceptual framework that encompasses a range of techniques—such as Rejection Sampling Fine-Tuning (RFT) \citep{yuan2023scaling}, Direct Preference Optimization (DPO) \citep{dpo}, PPO, and GRPO—interpreting them as direct or streamlined forms of reinforcement learning. Through a comprehensive set of experiments—spanning online versus offline training, result versus process supervision, single-step versus iterative reinforcement learning, among others—we probe the critical components underlying this framework. Finally, we elucidate the mechanisms by which reinforcement learning improves instruction-tuned model performance and outline avenues for advancing RL under this unified framework.

\subsection{Contributions}
This work presents advancements in scalable math pre-training and delves into the application and evaluation of reinforcement learning techniques.

\textbf{Math Pre-Training at Scale}

\begin{itemize}[topsep=0pt]
    \item We demonstrate that publicly available Common Crawl data harbors substantial mathematical information. Through a carefully engineered data selection process, we assemble the DeepSeekMath Corpus, a high-caliber dataset consisting of 120B tokens extracted from web content curated for mathematical relevance. This corpus is nearly 7 times larger than the set of math-related web pages employed by Minerva \citep{minerva}, and approximately 9 times the size of OpenWebMath \citep{openwebmath}.

    \item The DeepSeekMath-Base 7B pre-trained model achieves results on par with Minerva 540B \citep{minerva}, suggesting that parameter count alone does not determine mathematical reasoning strength. Effective pre-training on curated, high-quality mathematical data allows smaller models to attain competitive performance.

    \item Our experimental results shed light on the impact of pre-training order: code pre-training prior to mathematical training notably enhances problem-solving ability, both in cases involving tools and those that do not. This contributes partial resolution to the ongoing debate: \textit{does code pre-training support reasoning skill development?} Our evidence affirms its positive effect, at least with respect to mathematical reasoning.

    \item Contrary to common practice, incorporating arXiv paper data—while widespread in math-focused model pre-training—yields no discernible gains across the mathematical evaluation benchmarks assessed in this work.
\end{itemize}

\textbf{Exploration and Analysis of Reinforcement Learning}

\begin{itemize}[topsep=0pt]
    \item We present Group Relative Policy Optimization (GRPO), a resource-efficient reinforcement learning technique. By eliminating the need for a critic network and instead utilizing group-based baseline estimates, GRPO achieves substantial reductions in computational demands relative to Proximal Policy Optimization (PPO).
    \item We show that the application of GRPO leads to notable improvements in our instruction-tuned model, DeepSeekMath-Instruct, using only instruction-tuning datasets. Additionally, we find that reinforcement learning with GRPO contributes to better generalization to out-of-domain tasks.
    \item We offer a cohesive framework for interpreting various algorithms—including RFT, DPO, PPO, and GRPO—and systematically investigate their components through a broad set of experiments. These encompass online versus offline learning, comparisons of outcome and process-based supervision, single-turn against iterative reinforcement learning, and further dimensions, all aimed at elucidating the fundamental aspects underlying this framework.
    \item Guided by our unified perspective, we examine the factors responsible for reinforcement learning’s effectiveness and discuss prospective strategies for advancing reinforcement learning in large language models.
\end{itemize}

\subsection{Summary of Evaluations and Metrics}

\begin{itemize}[topsep=0pt]
    \item \textbf{English and Chinese Mathematical Reasoning}: Our models are thoroughly evaluated on mathematical reasoning tasks in both English and Chinese, spanning a range from elementary to university-level problems. The English evaluation suite comprises GSM8K \citep{gsm8k}, MATH \citep{MATH}, SAT \citep{llemma}, OCW Courses \citep{minerva}, and MMLU-STEM \citep{mmlu}; the Chinese evaluation suite includes MGSM-zh \citep{mgsm}, CMATH \citep{wei2023cmath}, Gaokao-MathCloze \citep{agieval}, and Gaokao-MathQA \citep{agieval}. Our assessment measures both the capability to produce independent, tool-free text solutions and to tackle problems programmatically via Python.

    On English datasets, DeepSeekMath-Base performs at a comparable level to the proprietary Minerva 540B \citep{minerva}, and outperforms all open-source foundation models (such as Mistral 7B \citep{mistral} and Llemma-34B \citep{llemma}), regardless of prior mathematical pre-training—often by a wide margin. In the context of Chinese benchmarks, DeepSeekMath-Base distinguishes itself as well, which we attribute to our incorporation of diverse, high-quality multilingual pre-training data rather than restricting to English as done in previous studies \citep{minerva,llemma}. With the introduction of mathematical instruction tuning and reinforcement learning, both DeepSeekMath-Instruct and DeepSeekMath-RL deliver impressive results, becoming the first open-source models to surpass the 50\% accuracy threshold on the challenging competition-level MATH dataset.
\end{itemize}

\item \textbf{Formal Mathematics}: To assess DeepSeekMath-Base, we utilize the informal-to-formal theorem proving benchmark introduced in \citep{dsp_proof}, conducted on the miniF2F dataset \citep{minif2f} with Isabelle \citep{isabelle} as the designated proof assistant. Our findings indicate that DeepSeekMath-Base achieves robust few-shot autoformalization results.

\item \textbf{Natural Language Understanding, Reasoning, and Code}: To comprehensively evaluate the models’ abilities in understanding, reasoning, and code generation, DeepSeekMath-Base is tested on several prominent benchmarks: Massive Multitask Language Understanding (MMLU) \citep{mmlu}, which includes 57 diverse multiple-choice tasks; BIG-Bench Hard (BBH) \citep{bbh}, a suite of 23 challenging tasks focused mainly on multi-step reasoning; along with HumanEval \citep{codex} and MBPP \citep{mbpp}, both standard datasets for assessing code-oriented language models. Math pre-training is found to enhance both linguistic reasoning and comprehension capabilities.

\section{Math Pre-Training}

\subsection{Data Collection and Decontamination}
This section details the methodology for developing the DeepSeekMath Corpus from Common Crawl. Figure \ref{fig:data_collect} illustrates the iterative framework employed to systematically extract a large-scale mathematical text corpus from Common Crawl, beginning with a high-quality seed dataset relevant to mathematics. Notably, this pipeline can be generalized to other specialized domains, including programming.

Our process starts by designating OpenWebMath \citep{openwebmath}—a curated set of high-caliber mathematical web documents—as the seed corpus. A fastText model \citep{joulin2016fasttext} is then trained to retrieve additional web content exhibiting characteristics similar to those in OpenWebMath. For this purpose, we randomly draw 500,000 samples from the seed as positive instances, pairing them with an equal number of negatives randomly selected from Common Crawl. Model training is performed using an open-source implementation\footnote{\url{https://fasttext.cc}} with the following configuration: vector dimension set to 256, learning rate at 0.1, word n-gram length capped at 3, minimum word frequency of 3, and 3 epochs. To effectively downsize the Common Crawl dataset, we first apply URL-based deduplication and near-duplicate removal, narrowing the data down to 40B HTML documents. The fastText model is then used to select mathematical web pages from this deduplicated dataset. Pages are scored by the fastText model, and only the highest-scoring ones are retained to filter out irrelevant or low-quality content. The retained data is further evaluated through pre-training trials across token thresholds of 40B, 80B, 120B, and 160B. In the initial cycle, we opt to retain the top 40B tokens.

Following the initial round of data acquisition, a significant number of mathematical web pages are still not retrieved. This shortfall is predominantly attributed to the insufficiently varied set of positive samples used in training the fastText model. To address this limitation, we expand the seed corpus by incorporating additional sources of mathematical content. The optimization process begins by segmenting the complete Common Crawl dataset into separate domains, where each domain encompasses web pages that share the same base URL. For each domain, we compute the fraction of web pages gathered during the first pass. Domains in which more than 10\% of pages have been harvested are designated as math-related (for example, \url{mathoverflow.net}). 

Next, we conduct a manual review of URLs within these selected domains, specifically annotating those that relate to mathematical content (such as \url{mathoverflow.net/questions}). Previously uncollected web pages connected to these URLs are incorporated into the seed corpus. This targeted inclusion enriches the collection of positive examples, thereby enhancing the performance of the fastText model in future data collection cycles by improving its ability to identify mathematical material. After four successive iterations, our collection comprises 35.5 million mathematical web pages, accumulating 120 billion tokens. Observing that nearly 98\% of this content was already collected in the third iteration, we decide to terminate further data collection after the fourth round.

In order to prevent data leakage into evaluation sets (benchmark contamination), we adopt the filtering methodology introduced by \cite{deepseek-coder}. We systematically remove web pages containing questions or answers drawn from established English mathematics benchmarks such as GSM8K \citep{gsm8k} and MATH \citep{MATH}, as well as Chinese benchmarks like CMATH \citep{wei2023cmath} and AGIEval \citep{agieval}. The removal process operates as follows: if any section of text contains a 10-gram sequence that exactly matches a substring from the evaluation benchmarks, the page is excluded from the training dataset. For evaluation texts that are under 10 grams but at least 3 grams in length, exact matching is applied to ensure that contaminated pages are effectively filtered.

\subsection{Validating the Quality of the DeepSeekMath~Corpus}

To assess the relative quality of the DeepSeekMath~Corpus, we conduct pre-training experiments comparing it against other recent mathematical datasets:

\begin{itemize}[topsep=0pt]
    \item \textbf{MathPile} \citep{mathpile}: A multi-source dataset of 8.9 billion tokens, comprising material from textbooks, Wikipedia, ProofWiki, CommonCrawl, StackExchange, and arXiv, with more than 85\% of content originating from arXiv;
    \item \textbf{OpenWebMath} \citep{openwebmath}: A subset of CommonCrawl refined for mathematical relevance, totaling 13.6 billion tokens;
    \item \textbf{Proof-Pile-2} \citep{llemma}: A mathematical dataset merging OpenWebMath, AlgebraicStack (with 10.3 billion tokens of mathematical code), and arXiv manuscripts (28.0 billion tokens). When conducting experiments with Proof-Pile-2, we adhere to the arXiv:Web:Code ratio of 2:4:1 as recommended in \cite{llemma}.
\end{itemize}

\subsubsection{Training Setting}
\label{sec:quality-policy}
Mathematical instruction tuning is conducted on a 1.3B parameter general-purpose language model that adopts the architecture used in DeepSeek LLMs \citep{deepseek-llm}, hereafter referred to as DeepSeek-LLM 1.3B. For each mathematical dataset, a separate model instance is fine-tuned over 150B tokens. All training runs utilize HAI-LLM \citep{haillm}, a resource-efficient and lightweight training platform. Consistent with DeepSeek LLMs' optimization procedures, the AdamW optimizer \citep{adamW} is utilized with parameters $\beta_1=0.9$, $\beta_2=0.95$, and a weight decay of $0.1$. The learning rate schedule employs multiple phases: the learning rate linearly warms up to its maximum value in 2,000 steps, then reduces to $31.6\%$ of its peak by $80\%$ progress through training, and further decays to $10.0\%$ at $90\%$ completion. The peak learning rate is set at $5.3\text{e-}4$, with batch sizes reaching 4M tokens and a maximum context window of 4K.

\subsubsection{Evaluation Results}

\textbf{The DeepSeekMath~Corpus stands out for its superior quality, broad multilingual coverage, and unparalleled scale.}

\begin{itemize}[topsep=0pt]
    \item \textbf{High-quality}:
    To assess downstream effectiveness, eight mathematical tasks are evaluated under few-shot chain-of-thought prompting \cite{cot}. Table \ref{tab:corpora_comparison} highlights the significant outperformance of models trained on DeepSeekMath~Corpus. According to Figure \ref{fig:corpora_comparisons}, this model exceeds the Proof-Pile-2-trained counterpart after processing 50B tokens (i.e., a complete pass through Proof-Pile-2), evidencing the higher per-token quality of DeepSeekMath~Corpus.
    \item \textbf{Multilingual}:
    Data within the DeepSeekMath~Corpus spans several languages, with English and Chinese as the dominant representations. Table \ref{tab:corpora_comparison} shows that leveraging DeepSeekMath~Corpus during training leads to notable gains in mathematical reasoning for both English and Chinese datasets. In comparison, earlier math corpora, with a heavy focus on English, either show minimal effects or actively degrade Chinese mathematical reasoning performance.
    \item \textbf{Large-scale}:
    The DeepSeekMath~Corpus is several orders of magnitude larger than other existing mathematical datasets. Figure \ref{fig:corpora_comparisons} illustrates that DeepSeek-LLM 1.3B, when trained on DeepSeekMath~Corpus, achieves a more rapid learning trajectory and maintains sustained improvements over time. Conversely, models trained with smaller baseline corpora quickly exhaust novel data after multiple epochs, causing model progress to stagnate.
\end{itemize}

\subsection{Training and Evaluating DeepSeekMath-Base 7B}

Here, we present DeepSeekMath-Base 7B, a foundational model tailored for advanced mathematical reasoning tasks. This model is initialized using DeepSeek-Coder-Base-v1.5 7B \citep{deepseek-coder} and subsequently trained on a corpus of 500B tokens. The composition of the training data includes 56\% from the DeepSeekMath Corpus, 4\% derived from AlgebraicStack, 10\% from arXiv, 20\% sourced from Github code repositories, and the remaining 10\% from Common Crawl, featuring natural language content in both English and Chinese. Training settings follow those outlined in Section \ref{sec:quality-policy}, with two primary deviations: the learning rate peaks at 4.2e-4, and each batch encompasses 10M tokens.

A thorough evaluation is performed on DeepSeekMath-Base 7B to measure its proficiency in mathematics, emphasizing three core areas: generating standalone mathematical solutions independent of auxiliary tools, addressing mathematical tasks via tool augmentation, and engaging in rigorous formal theorem proving. The model's general competencies—such as natural language understanding, logical reasoning, and coding ability—are also benchmarked for a comprehensive performance overview.

\paragraph{Mathematical Problem Solving with Step-by-Step Reasoning}

DeepSeekMath-Base is assessed on its capacity to solve mathematical questions using a few-shot chain-of-thought prompting paradigm \citep{cot} across eight benchmark datasets in both English and Chinese. The selected benchmarks span a spectrum of quantitative reasoning challenges (including GSM8K \citep{gsm8k}, MATH \citep{MATH}, and CMATH \citep{wei2023cmath}) as well as multiple-choice question sets (such as MMLU-STEM \citep{mmlu} and Gaokao-MathQA \citep{agieval}), collectively representing a broad range of mathematical disciplines and difficulty levels, from basic to advanced collegiate problems.

As detailed in Table \ref{tab:base_math_cot}, DeepSeekMath-Base 7B consistently achieves the highest scores across all eight benchmarks when compared to other open-source base models, such as the prevalent Mistral 7B \citep{mistral} and the newly introduced Llemma 34B \citep{llemma}, which is specifically math-trained on Proof-Pile-2 \citep{llemma}. Of particular significance, DeepSeekMath-Base demonstrates a more than 10\% absolute lead over prior open-source base models on the advanced MATH competition dataset, even outperforming Minerva 540B \citep{minerva}—a proprietary model that is 77 times larger, constructed upon PaLM \citep{palm}, and further refined on extensive mathematical corpora.

\paragraph{Mathematical Problem Solving with Tool Use} 
We assess models’ ability to perform program-assisted mathematical reasoning using GSM8K and MATH datasets under the few-shot program-of-thought prompting framework \citep{pot,pal}. In this setup, the models are guided to solve problems through generating Python code, utilizing libraries such as \textit{math} and \textit{sympy} for advanced calculations. The correctness of answers is determined by executing these programs. According to the results in Table \ref{tab:base_math_pot_proof}, DeepSeekMath-Base 7B achieves superior results compared to the previous leading open-source model, Llemma 34B.

\paragraph{Formal Mathematics}

Automating formal proofs enhances the rigor and dependability of mathematical verification, as well as overall workflow efficiency, and has attracted growing research focus. We evaluate DeepSeekMath-Base 7B using the informal-to-formal proof generation task from \citep{dsp_proof}, where the model converts an informal statement, along with its formal representation and informal proof, into a formal proof. Using the miniF2F benchmark \citep{minif2f}, which targets Olympiad-grade formal mathematics, the model is prompted in a few-shot manner to produce formal proofs in Isabelle for each problem. Mirroring the methodology in \cite{dsp_proof}, we utilize the model to draft proof outlines, then apply Sledgehammer \citep{sledgehammer}, a standard automated theorem prover, to complete any omitted details. Table \ref{tab:base_math_pot_proof} indicates that DeepSeekMath-Base 7B displays notable proficiency in automating the formalization of proofs.

\paragraph{Natural Language Understanding, Reasoning, and Code}

To gauge language comprehension, reasoning, and coding skills, we benchmark performance on MMLU \citep{mmlu} (natural language understanding), BBH \citep{bbh} (reasoning), HumanEval \citep{codex}, and MBPP \citep{mbpp} (coding tasks). As shown in Table \ref{tab:understanding_reasoning_code}, DeepSeekMath-Base 7B shows considerable improvement over its predecessor, DeepSeek-Coder-Base-v1.5 \citep{deepseek-coder}, particularly in MMLU and BBH, highlighting the advantage of targeted mathematical training for general language and reasoning tasks. Furthermore, by incorporating code tokens into its ongoing training regimen, DeepSeekMath-Base 7B manages to preserve the strong coding performance characteristic of DeepSeek-Coder-Base-v1.5 on both code-related benchmarks. On the whole, DeepSeekMath-Base 7B clearly outperforms the general-purpose Mistral 7B \citep{mistral} on all three reasoning and programming evaluation tasks.

\section{Supervised Fine-Tuning}

\subsection{SFT Data Curation}

We assemble a dataset for mathematical instruction-tuning that spans both English and Chinese problem sets, drawing from a range of mathematical disciplines and complexity levels. Each problem is paired with corresponding solutions using chain-of-thought (CoT) \citep{cot}, program-of-thought (PoT) \citep{pot,pal}, as well as tool-integrated reasoning approaches \citep{tora}. The total dataset comprises 776K training examples.

\begin{itemize}[topsep=0pt]
    \item \textbf{English mathematical datasets}: We provide tool-integrated solution annotations for GSM8K and MATH datasets, incorporate a selected subset from MathInstruct \citep{MathInstruct}, and utilize the training data from Lila-OOD \citep{lila}, where solutions follow either CoT or PoT methodologies. This English dataset encompasses a broad array of mathematical domains, such as algebra, probability, number theory, calculus, and geometry.
    \item \textbf{Chinese mathematical datasets}: For the Chinese component, we curate K-12 level math problems from 76 sub-domains—including topics like linear equations—with solutions annotated using both CoT and tool-integrated reasoning.
\end{itemize}

\subsection{Training and Evaluating DeepSeekMath-Instruct 7B}

This section details the supervised fine-tuning of DeepSeekMath-Instruct 7B, which is based on the DeepSeekMath-Base model. Training samples are concatenated at random until a maximum context length of 4K tokens is reached. The model undergoes 500 training steps using a batch size of 256, with the learning rate set at 5e-5.

Model evaluation is carried out in terms of mathematical reasoning performance, both in the absence and presence of tool support, across four quantitative reasoning benchmarks in both English and Chinese. We compare our model’s outcomes against top-performing models available at the time:

\begin{itemize}[topsep=0pt]
    \item \textbf{Closed-source models}: These include (1) the GPT series, with GPT-4 \citep{gpt4} and GPT-4 Code Interpreter~\footnote{\url{https://openai.com/blog/chatgpt-plugins\#\#code-interpreter}} representing the highest-performing variants, (2) Gemini Ultra and Pro \citep{gemini}, (3) Inflection-2 \citep{inflection-2}, (4) Grok-1~\footnote{\url{https://x.ai/model-card}}, as well as recently introduced Chinese models such as (5) Baichuan-3~\footnote{\url{https://www.baichuan-ai.com}}, and (6) the latest version of GLM-4~\footnote{\url{https://open.bigmodel.cn/dev/api\#glm-4}} from the GLM series \citep{glm}.
\end{itemize}

These models serve broad, general applications, and most have passed through multiple alignment stages.
    \item \textbf{Open-source models} consist of: general-purpose models such as (1) DeepSeek-LLM-Chat 67B \citep{deepseek-llm}, (2) Qwen 72B \citep{qwen}, (3) SeaLLM-v2 7B \citep{seallm}, and (4) ChatGLM3 6B \citep{chatglm3}. There are also models with specific improvements in mathematical reasoning: (5) InternLM2-Math 20B~\footnote{\url{https://github.com/InternLM/InternLM-Math}}, which extends InternLM2 by incorporating math-specific training before further instruction tuning; (6) Math-Shepherd-Mistral 7B, which utilizes PPO training \citep{schulman2017proximal} on Mistral 7B \citep{mistral} with a process-supervised reward model; (7) the WizardMath family \citep{wizardmath}, which boosts mathematical reasoning for Mistral 7B and Llama-2 70B \citep{llama2} via evolve-instruct (an instruction tuning variant with AI-generated prompts) and PPO training, predominantly leveraging problems from GSM8K and MATH; (8) MetaMath 70B \citep{metamath}, produced by fine-tuning Llama-2 70B on an expanded version of GSM8K and MATH; (9) ToRA 34B \cite{tora}, which is CodeLlama 34B optimized for tool-augmented mathematical reasoning; and (10) MAmmoTH 70B \citep{MathInstruct}, created by instruction-tuning Llama-2 70B on MathInstruct.
\end{itemize}

As indicated in Table \ref{tab:sft_rl_math}, when tool usage is prohibited during evaluation, DeepSeekMath-Instruct 7B displays robust performance for stepwise reasoning. Remarkably, on the competition-grade MATH dataset, this model exceeds all open-source alternatives and outperforms most closed-source models (e.g., Inflection-2, Gemini Pro) by at least 9\% in absolute terms. This advantage is maintained even when compared against much larger models (e.g., Qwen 72B) or those trained with advanced math-oriented reinforcement learning techniques (e.g., WizardMath-v1.1 7B). Although DeepSeekMath-Instruct matches the Chinese proprietary models GLM-4 and Baichuan-3 on the MATH benchmark, it still lags behind top-tier systems like GPT-4 and Gemini Ultra.

Allowing models to combine natural language reasoning with code-based tool use during evaluation, DeepSeekMath-Instruct 7B achieves nearly 60\% accuracy on MATH, setting a new high-water mark among open-source models. Across additional benchmarks, its results are comparable to DeepSeek-LLM-Chat 67B, the former state-of-the-art that boasts ten times the parameters.

\section{Reinforcement Learning}

\subsection{Group Relative Policy Optimization} Reinforcement learning (RL) has shown notable effectiveness in further enhancing the mathematical reasoning skills of LLMs following the Supervised Fine-Tuning (SFT) phase \citep{wang2023math,wizardmath}. Here, we present Group Relative Policy Optimization (GRPO), our resource-efficient and high-performance RL technique.

\subsubsection{From PPO to GRPO} Proximal Policy Optimization (PPO) \citep{schulman2017proximal} represents an actor-critic algorithm that is frequently applied in the RL fine-tuning phase for LLMs \citep{ouyang2022training}. Its goal is to optimize the LLM policy through the maximization of the following surrogate objective:

\begin{equation}
\footnotesize
    \mathcal{J}_{PPO}(\theta) = \mathbb{E}{[q \sim P(Q), o \sim \pi_{\theta_{old}}(O|q)]} \frac{1}{|o|} \sum_{t=1}^{|o|} \min \left[ \frac{\pi_\theta(o_{t} | q, o_{<t})}{\pi_{\theta_{old}}(o_{t} | q, o_{<t})} A_{t}, \text{clip} \left( \frac{\pi_\theta(o_{t} | q, o_{<t})}{\pi_{\theta_{old}}(o_{t} | q, o_{<t})}, 1 - \varepsilon,

Here, $\pi_{\theta}$ and $\pi_{\theta_{old}}$ denote the current and the previous policy models, respectively, and $q, o$ correspond to questions and outputs that are drawn from the question dataset and generated using the old policy $\pi_{\theta_{old}}$. The hyper-parameter $\varepsilon$ is used in Proximal Policy Optimization (PPO) as a clipping parameter to stabilize learning. The term $A_t$ represents the advantage, which is estimated by applying Generalized Advantage Estimation (GAE) \citep{gae} utilizing future rewards $\{r_{\ge t}\}$ together with a learned value function $V_{\psi}$. Therefore, PPO requires joint training of a value function alongside the policy, and to guard against excessive optimization toward the reward model, it is standard practice to include a per-token KL penalty based on a reference model in the reward at every token \citep{ouyang2022training}, formalized as

\begin{equation}
    r_{t} = r_\varphi(q, o_{\le t}) - \beta \log\frac{\pi_{\theta}(o_{t}|q, o_{<t})}{\pi_{ref}(o_{t}|q, o_{<t})},
\label{eq:PPO-reward}
\end{equation}

where $r_\varphi$ is the reward function, $\pi_{ref}$ refers to a reference policy—commonly set as the starting SFT model—and $\beta$ controls the magnitude of the KL penalty.

Given that the value function employed by PPO is often of comparable scale to the policy model, this introduces notable overheads in both computation and memory. Furthermore, PPO relies on the value function as a baseline to reduce variance in advantage estimates during RL training. In large language model (LLM) settings, the reward model typically only assigns a score for the final token of the generated sequence, posing challenges for accurate value function training at every timestep. To overcome these limitations, as illustrated in Figure \ref{fig:grpo}, we introduce Group Relative Policy Optimization (GRPO), which eliminates the need for a separate value function as in PPO. Instead, GRPO leverages the average reward across multiple sampled responses to a given question as a baseline. Specifically, for each question $q$, a collection of outputs $\{o_1, o_2, \cdots, o_G\}$ is sampled from the old policy $\pi_{\theta_{old}}$, and the policy is updated by maximizing

\begin{equation}
\footnotesize
\begin{split}
    \mathcal{J}_{GRPO}(\theta) &= \mathbb{E}{[q \sim P(Q), \{o_i\}_{i=1}^G \sim \pi_{\theta_{old}}(O|q)]}  \\
    & \frac{1}{G}\sum_{i=1}^G\frac{1}{|o_i|} \sum_{t=1}^{|o_i|} \left\{ \min \left[ \frac{\pi_\theta(o_{i,t} | q, o_{i,<t})}{\pi_{\theta_{old}}(o_{i,t} | q, o_{i,<t})} \hat{A}_{i,t}, \text{clip} \left( \frac{\pi_\theta(o_{i,t} | q, o_{i,<t})}{\pi_{\theta_{old}}(o_{i,t} | q, o_{i,<t})}, 1 - \varepsilon, 1 + \varepsilon \right)  \hat{A}_{i,t} \right] - \beta \mathbb{D}_{KL}\left[\pi_{\theta} || \pi_{ref}\right]\right\} ,
\end{split}
\label{eq:GRPO-obj}
\end{equation}

where the hyper-parameters $\varepsilon$ and $\beta$ serve the same purpose as before, and $\hat{A}_{i,t}$ is the advantage derived using only relative rewards from the outputs in the corresponding group, as will be further discussed in subsequent sections. This group-relative computation of the advantage in GRPO naturally fits with the comparative framework in which reward models are usually trained, as those models rely on pairwise output comparisons for a single prompt. It is also noteworthy that, in contrast to incorporating the KL penalty directly in the reward as done in PPO, GRPO enforces regularization by including the KL divergence between the learned policy and the reference policy as an explicit penalty in the objective, thereby simplifying the computation of $\hat{A}_{i,t}$. Furthermore, unlike the KL term present in (\ref{eq:PPO-reward}),

\begin{algorithm}[t]
  \small
  \caption{Iterative Group Relative Policy Optimization}
  \textbf{Input} initial policy model $\pi_{\theta_{\text{init}}}$; reward models $r_\varphi$; task prompts $\mathcal{D}$; 
  hyperparameters $\varepsilon$, $\beta$, $\mu$
  \begin{algorithmic}[1]
    \State Initialize policy model: $\pi_\theta \leftarrow \pi_{\theta_{\text{init}}}$
    \For{iteration = 1, \dots, I}
       \State Set reference model $\pi_{ref} \leftarrow \pi_{\theta}$
      \For{step = 1, \dots, M}
      \State Draw a mini-batch $\mathcal{D}_b$ from $\mathcal{D}$
      \State Assign old policy $\pi_{\theta_{old}} \leftarrow \pi_{\theta}$ 
      \State For every question $q \in \mathcal{D}_b$, generate $G$ samples $\{o_i\}_{i=1}^G \sim \pi_{\theta_{old}} (\cdot \mid q)$
      \State Calculate corresponding rewards $\{r_i\}_{i=1}^G$ for each output $o_i$ using $r_{\varphi}$
      \State For each token $t$ in $o_i$, compute $\hat{A}_{i,t}$ via group relative advantage estimation
      \For{GRPO iteration = 1, \dots, $\mu$}
        \State Update $\pi_{\theta}$ to increase the GRPO objective (Equation \ref{eq:GC-GRPO})
      \EndFor
    \EndFor \State Continue training $r_\varphi$ with replay data for ongoing learning.
    \EndFor 
  \end{algorithmic}
  \textbf{Output} $\pi_\theta$
  \label{alg:iter-grpo}
\end{algorithm}

\subsubsection{Outcome Supervision RL with GRPO} Given a prompt $q$, a collection of responses $\{o_1, o_2, \cdots, o_G\}$ is sampled from the prior policy model $\pi_{\theta_{old}}$. A reward model assigns a scalar reward to each response, resulting in a set of $G$ rewards $\mathbf{r} = \{r_1, r_2, \cdots, r_G\}$. These raw rewards undergo normalization by removing the mean of the group and dividing by their standard deviation. For outcome supervision, this normalized score is used as the reward for the entire output $o_i$, and all tokens in $o_i$ share the same advantage estimate: $\hat{A}_{i, t} = \widetilde{r}_i = \frac{r_i- {\rm mean}(\mathbf{r})}{{\rm std}(\mathbf{r})}$. The policy update proceeds by optimizing the objective laid out in equation (\ref{eq:GRPO-obj}).

\subsubsection{Process Supervision RL with GRPO} While outcome supervision provides a single terminal reward per output, it often falls short for complex mathematical reasoning tasks that benefit from more granular feedback. Following the approach of \cite{wang2023math}, we also adopt process supervision, wherein rewards are allocated after each reasoning step. Concretely, given $q$ and the batch of $G$ generated outputs $\{o_1, o_2, \cdots, o_G\}$, a process reward model computes per-step rewards, generating
$\mathbf{R} = \{ \{r_1^{index(1)},\cdots,r_1^{index(K_1)}\}, \cdots,  \{r_G^{index(1)},\cdots,r_G^{index(K_G)}\} \}$,
where $index(j)$ is the final token position of the $j$-th reasoning step in $o_i$, and $K_i$ is the number of steps in output $i$. The set of step-wise rewards is normalized across all steps and all outputs as $\widetilde{r}_i^{index(j)} = \frac{r_i^{index(j)} - {\rm mean(\mathbf{R})}}{{\rm std(\mathbf{R})}}$.
The advantage estimate for each token $t$ in $o_i$ is given by summing all normalized future step rewards: $\hat{A}_{i, t} = \sum_{index(j) \ge t} \widetilde{r}_i^{index(j)}$.
The policy is then refined by maximizing the objective described in equation (\ref{eq:

\subsubsection{Iterative RL with GRPO} As RL training advances, the previously trained reward model may become inadequate for effectively guiding the evolving policy model. To address this, we implement an iterative RL procedure utilizing GRPO. In this approach, detailed in Algorithm \ref{alg:iter-grpo}, we continually augment the reward model’s training data with samples produced by the updated policy model, employing a replay strategy that retains 10\% of prior data. Subsequently, the current policy model is assigned as the reference model, and policy model updates proceed using the most recent reward model.

\subsection{Training and Evaluating DeepSeekMath-RL}

Our RL experiments are based on the DeepSeekMath-Instruct 7B model. The RL phase utilizes questions in the chain-of-thought format drawn from GSM8K and MATH subsets within the SFT corpus, totaling approximately 144K examples. We intentionally omit other SFT-derived questions to isolate the effect of RL on benchmarks that do not have training data present during RL. Reward model training datasets are constructed as per the methodology of \citep{wang2023math}. The initial reward model is trained from DeepSeekMath-Base 7B, applying a learning rate of 2e-5. For GRPO optimization, the policy model adopts a learning rate of 1e-6 and a KL coefficient of 0.04. Each prompt is used to sample $64$ outputs, with a maximum sequence length of 1024 tokens and a batch size of 1024. The policy model undergoes one update at the conclusion of every exploration phase. Performance evaluation of DeepSeekMath-RL 7B is carried out on the same benchmarks as DeepSeekMath-Instruct 7B. For the DeepSeekMath-RL 7B model, the tasks involving GSM8K and MATH with chain-of-thought reasoning are categorized as in-domain, whereas all remaining benchmarks are treated as out-of-domain.

Table \ref{tab:sft_rl_math} presents a comparison between open- and closed-source models leveraging chain-of-thought and tool-augmented reasoning strategies on both English and Chinese benchmarks. The results indicate the following: (1) DeepSeekMath-RL 7B achieves chain-of-thought reasoning accuracies of 88.2\% on GSM8K and 51.7\% on MATH, outperforming every open-source model in the 7B to 70B parameter range and exceeding the results of most closed-source models as well. (2) Importantly, DeepSeekMath-RL 7B is exclusively fine-tuned on chain-of-thought-style instruction data from GSM8K and MATH, with DeepSeekMath-Instruct 7B as the initialization point. Despite its narrowly focused training data, it achieves superior performance to DeepSeekMath-Instruct 7B on all metrics, highlighting the value added by reinforcement learning.

\section{Discussion}
This section discusses insights gathered from our pre-training and reinforcement learning experiments.

\subsection{Lessons Learnt in Pre-Training}

We begin with reflections on the pre-training stage. Unless otherwise stated, we maintain the training configuration specified in Section \ref{sec:quality-policy}. Notably, the DeepSeekMath Corpus referenced here consists of an 89B-token dataset produced during the second data collection round.

\subsubsection{Code Training Enhances Mathematical Reasoning}

It has been widely conjectured, though not systematically proven, that incorporating code in training regimens benefits reasoning skills. Here, we provide supporting evidence within the mathematical context: models trained with code demonstrate improved mathematical reasoning capabilities, both when using and not using external tools.

To investigate the impact of code training on mathematical reasoning, we carried out comparisons between models trained with two distinct approaches: two-stage training and one-stage training.

\textbf{Two-Stage Training}

\begin{itemize}[topsep=0pt]
    \item \textbf{Code Training for 400B Tokens $\rightarrow$ Math Training for 150B Tokens}: DeepSeek-LLM 1.3B is initially trained on 400B code tokens, after which it undergoes an additional 150B math token training phase;
    \item \textbf{General Training for 400B Tokens $\rightarrow$ Math Training for 150B Tokens}: To serve as a control, we substitute the code token pretraining with 400B general tokens, sampled from DeepSeek-AI's extensive general corpus. This design aims to determine if code tokens offer advantages over general-domain tokens in facilitating improved mathematical reasoning capabilities.
\end{itemize}

\textbf{One-Stage Training}

\begin{itemize}[topsep=0pt]
    \item \textbf{Math Training for 150B Tokens}: We expose DeepSeek-LLM 1.3B solely to 150B math tokens for training;
    \item \textbf{Training on a mixture of 400B Code Tokens and 150B Math Tokens}: Given that sequential math finetuning after code pretraining has been observed to reduce performance on coding tasks, we also consider an alternative in which math and code tokens are combined in a single-stage regimen. This approach is intended to assess whether mathematical reasoning can still be improved while concurrently alleviating catastrophic forgetting of coding knowledge.
\end{itemize}

\paragraph{Results}
The performance outcomes corresponding to various training methodologies are summarized in Table \ref{tab:code-math} and Table \ref{tab:code-math-general-eval}.

Pretraining on code notably enhances performance in program-assisted mathematical reasoning tasks, effective across both sequential (two-stage) and integrated (one-stage) training paradigms. Table \ref{tab:code-math} reveals that code pretraining alone robustly strengthens DeepSeek-LLM’s ability to address GSM8K and MATH datasets via Python. Incorporating an additional math training phase leads to even higher performance. Moreover, the one-stage combined training strategy—simultaneously mixing code and math tokens—successfully mitigates catastrophic forgetting (a typical pitfall of sequential approaches), while concurrently improving both code generation (Table \ref{tab:code-math-general-eval}) and program-mediated mathematical reasoning (Table \ref{tab:code-math}).

Furthermore, code-centric pretraining yields benefits for mathematical reasoning tasks even when no external tools are employed. With two-stage training, the model already demonstrates moderate gains following the initial code phase and achieves maximal efficacy after subsequent math finetuning. However, in the joint training setting where code and math tokens are presented together from the outset, performance on tool-free mathematical reasoning tasks declines. This outcome may be attributed to DeepSeek-LLM 1.3B's constrained capacity, which could prevent the simultaneous internalization of both code and mathematics at this scale.

\subsubsection{ArXiv Papers Appear Unproductive for Enhancing Mathematical Reasoning}

While arXiv papers are frequently incorporated into pre-training datasets for mathematical models \citep{minerva,gpt-f,llemma,mathpile}, there has been limited rigorous investigation into their concrete contribution to mathematical reasoning abilities. Surprisingly, our findings suggest that arXiv papers do not significantly aid in fostering mathematical reasoning skills. We conducted experiments across a range of model sizes, notably DeepSeek-LLM 1.3B and DeepSeek-Coder-Base-v1.5 7B \citep{deepseek-coder}, utilizing two distinct arXiv-based corpora, each processed via different pipelines:

\begin{itemize}[topsep=0pt]
    \item \textbf{MathPile} \citep{mathpile}: an 8.9 billion token dataset curated using specialized cleaning and filtering strategies, with more than 85\% sourced from scientific arXiv submissions;
    \item \textbf{ArXiv-RedPajama} \citep{redpajama}: a dataset comprising all arXiv LaTeX files, systematically stripped of preambles, comments, macro definitions, and reference lists, totaling 28.0 billion tokens.
\end{itemize}

For our study, DeepSeek-LLM 1.3B was pre-trained for 150B tokens and DeepSeek-Coder-Base-v1.5 7B for 40B tokens on each of the arXiv corpora separately. The results indicate that exclusive training on arXiv data does not facilitate gains in mathematical reasoning; in some cases, model performance even regresses on a spectrum of mathematical evaluation sets considered. The evaluated tasks span quantitative reasoning datasets such as GSM8K and MATH (Table \ref{tab:arxiv-cot}), multiple-choice STEM assessments like MMLU-STEM (Table \ref{tab:arxiv-cot}), and formal mathematics datasets like miniF2F (Table \ref{tab:arxiv_atp}).

Nevertheless, this outcome is not conclusive, and several caveats apply. Our exploration did not address:

\begin{itemize}[topsep=0pt]
    \item The role of arXiv-sourced data in specialized mathematical tasks beyond our benchmark suite, such as converting formal theorems and proofs into informal language;
    \item How arXiv tokens perform when supplemented by other dataset modalities;
    \item Whether larger-scale models might unlock latent advantages of arXiv-based pre-training.
\end{itemize}

Consequently, more detailed research is warranted, and we defer these questions to future work.

\subsection{Reinforcement Learning Insights}
\subsubsection{Toward a Unified Training Framework} This section develops a general framework to systematically analyze diverse training paradigms—including SFT, RFT, DPO, PPO, and GRPO—supplemented by empirical investigations into key aspects of this unified approach. In broad terms, the training gradient with respect to the parameters $\theta$ for these methods may be represented as:

\begin{equation}
    \nabla_{\theta}\mathcal{J}_{\textcolor{red}{\mathcal{A}}}(\theta) = \mathbb{E}[\underbrace{(q,o) \sim \textcolor{red}{\mathcal{D}}}_{Data \ Source}]\left( \frac{1}{|o|} \sum_{t=1}^{|o|}  \underbrace{GC_{{\mathcal{A}}}(q, o, t, \textcolor{red}{\pi_{{rf}}})}_{Gradient \ Coefficient}  \nabla_{\theta}\log \pi_{\theta}(o_t | q, o_{<t})\right).
\label{eq:objective}
\end{equation}

This objective can be characterized by three principal elements: (1) \textit{Data Source $\mathcal{D}$}, which specifies where the training examples are drawn from; (2) \textit{Reward Function $\pi_{{rf}}$}, which serves as the basis for computing the reward signal during training; (3) \textit{Algorithm $\mathcal{A}$}, which integrates the data and reward to compute the gradient coefficient $GC$ that modulates the strength and direction of learning updates, either punishing or rewarding particular data points. Using this general formulation, several well-known training approaches can be systematically described:

\begin{itemize}
    \item  \textbf{Supervised Fine-tuning (SFT)}: SFT further trains a pretrained model by exposing it to a curated set of examples chosen by humans.
    \item \textbf{Rejection Sampling Fine-tuning (RFT)}: RFT enhances the SFT model by additionally training it on a subset of model-generated outputs, filtered for correctness, in response to prompts from the SFT data.
    \item \textbf{Direct Preference Optimization (DPO)}: DPO updates the SFT model by leveraging augmented response pairs produced from the SFT model, optimizing according to a pairwise DPO objective.
    \item \textbf{Online Rejection Sampling Fine-tuning (Online RFT)}: In contrast to standard RFT, Online RFT initializes with the SFT model but continually samples and filters outputs from the current version of the policy model for additional fine-tuning.
    \item \textbf{PPO/GRPO}: PPO/GRPO starts from the SFT-trained model as the policy and iteratively strengthens it by rewarding sampled responses drawn in real time from the evolving policy model.
\end{itemize}

We present a summary of the constituent elements of these approaches in Table \ref{tab:data_policy}. For a comprehensive discussion and step-by-step derivations, please consult Appendix \ref{app:analysis-rl}.

\paragraph{Insight on Data Source} We classify data sources into two main types: online sampling and offline sampling. Online sampling involves collecting training data from the ongoing exploration carried out by the policy model during training. In contrast, offline sampling utilizes training data derived from samples generated by the initial SFT model. The RFT and DPO methods are representative of the offline approach, whereas Online RFT and GRPO utilize the online approach.

Referencing Figure \ref{fig:rl-analysis}, our findings indicate that Online RFT markedly surpasses RFT on both evaluated benchmarks. Initially, both methods perform similarly, but as training progresses, Online RFT establishes a distinct performance lead, showcasing the benefits of online learning. This result is expected because at the outset, the actor closely mirrors the SFT model, and differences in the sampled data are minimal. As training advances, actor-generated samples begin to diverge more significantly, and leveraging real-time data becomes increasingly beneficial.

\paragraph{Insight on Gradient Coefficient} The reward signal is transformed into a gradient coefficient to guide updates to the model parameters. In our analysis, we categorize reward functions as either `Rule'-based or `Model'-based. The `Rule' approach assesses response quality by directly evaluating the correctness of the answer, while the `Model' approach involves training a reward model to assign scores to responses, with its training data informed by rule-based judgments. As demonstrated in Equations \ref{eq:GC-RFT} and \ref{eq:GC-GRPO}, a critical distinction between GRPO and Online RFT is that GRPO modulates its gradient coefficient according to the magnitude of rewards issued by the reward model, thereby enabling varying degrees of reinforcement or penalization based on response quality. Conversely, Online RFT does not discriminate between correct responses; it applies the same level of reinforcement for all answers deemed correct and does not apply negative updates to incorrect answers.

As illustrated in Figure \ref{fig:rl-analysis}, GRPO demonstrates superior performance compared to online RFT, underscoring the effectiveness of modulating the coefficients for positive and negative gradients. Moreover, GRPO+PS outperforms GRPO+OS, which suggests that employing finer-grained, step-dependent gradient coefficients confers additional advantages. We also examine iterative RL in our studies, where we implement two consecutive rounds of iteration. The results presented in Figure \ref{fig:iter-rl} reveal that applying iterative RL leads to notable gains, with the most pronounced improvement occurring after the initial iteration.

\subsubsection{Why RL Works?}  
This study employs reinforcement learning on a subset of the instruction tuning dataset, which results in marked gains over the baseline instruction tuning model. To clarify the underlying mechanisms of these improvements, we assess the Pass@K and Maj@K metrics for both the Instruct and RL-enhanced models across two benchmark datasets. As indicated in Figure \ref{fig:maj-pass}, reinforcement learning boosts the Maj@K metric but does not lead to improvements in Pass@K. This pattern implies that RL strengthens the robustness of the output distribution, effectively enhancing the likelihood of retrieving a correct answer from the TopK responses, rather than fundamentally advancing the model’s baseline ability. In a similar vein, \citep{wang2023making} identified a \textbf{misalignment problem} in reasoning challenges for SFT models, and demonstrated that SFT models' reasoning skills can be elevated through targeted preference alignment strategies \citep{yuan2023rrhf,song2023preference,wang2023making}.

\subsubsection{How to Achieve More Effective RL?}
\label{sec:effec_rl}
Our experiments indicate that RL is highly effective in addressing mathematical reasoning tasks. Furthermore, we introduce a unified perspective for interpreting several key training paradigms, framing all considered methods as either instances of direct RL or simplified RL variants. As encapsulated in Equation \ref{eq:objective}, these approaches involve three critical components: Data Source, Algorithm, and Reward Function. We suggest possible avenues for future research in relation to each component.

\paragraph{Data Source}  
The data source forms the foundation for any training procedure. Within RL, we interpret the data source as comprising unlabeled problems and outputs generated from the current policy model. In our present experiments, we limit ourselves to instruction tuning questions, paired with outputs obtained via basic nucleus sampling. We hypothesize that this restriction might contribute to our observation that RL mainly enhances Maj@K, rather than broader performance metrics. Future work will involve extending our RL framework to handle out-of-distribution prompts and implementing \textbf{advanced sampling (decoding) strategies}, such as those informed by tree-search algorithms \citep{yao2023tree}. Additionally, advances in \textbf{efficient inference techniques} \citep{xia-etal-2023-speculative,leviathan2023fast,kwon2023efficient,xia2024unlocking}, which directly affect the exploration capability of policy models, are anticipated to play a pivotal role.

\paragraph{Algorithms} Algorithms utilize the combination of data and reward signals to determine the gradient coefficient, which in turn guides updates to model parameters. According to Equation \ref{eq:objective}, virtually all contemporary approaches fundamentally \textbf{TRUST} the reward signal when adjusting the conditional probabilities of specific tokens, either increasing or decreasing them. Nevertheless, there is no guarantee that the reward signal will always be dependable, especially for highly intricate tasks. For instance, even in meticulously curated datasets like PRM800K \citep{lightman2023let}, annotated by skilled professionals, roughly 20\% of labels remain erroneous\footnote{\url{https://github.com/openai/prm800k/issues/12\#issuecomment-1728491852}}. Consequently, we propose examining reinforcement learning algorithms that exhibit robustness in the face of noisy or imperfect reward feedback. Our hypothesis is that \textbf{WEAK-TO-STRONG} alignment techniques \citep{burns2023weak} could fundamentally transform existing learning algorithm paradigms.

\paragraph{Reward Function} The reward function provides the primary learning signal in RL, typically represented by a neural reward model. We identify three central challenges for advancing reward models: 1) \textbf{Improving the reward model's ability to generalize.} To be effective, reward models must successfully extrapolate to novel, out-of-distribution queries and sophisticated decoding results; if not, RL may simply regularize the current output distribution of LLMs rather than enhancing their core performance; 2) \textbf{Incorporating uncertainty in reward estimation.} Quantifying uncertainty could serve as a critical interface between underpowered reward models and the aforementioned weak-to-strong training algorithms; 3) \textbf{Developing process reward models that are both high-quality and resource-efficient,} which are capable of delivering nuanced, detailed supervision for step-wise reasoning tasks \citep{lightman2023let,wang2023math}.

\section{Conclusion, Limitation, and Future Work} This work introduces DeepSeekMath, a model that establishes new state-of-the-art performance among open-source systems on the competition-level MATH benchmark and nearly matches the achievements of closed-source alternatives. DeepSeekMath~builds upon the DeepSeek-Coder-v1.5 7B foundation and is continually trained for 500B tokens, with 120B math-specific tokens derived primarily from Common Crawl forming a substantial segment of the training data. Through detailed ablation analysis, we demonstrate that web-sourced data offers remarkable promise for compiling high-quality mathematical resources, whereas arXiv-based content may not be as advantageous as initially anticipated. We further propose Group Relative Policy Optimization (GRPO), an adaptation of Proximal Policy Optimization (PPO), capable of significantly enhancing mathematical reasoning proficiency while reducing memory requirements. Experimental findings affirm that GRPO continues to deliver gains even when DeepSeekMath-Instruct 7B achieves strong benchmark results. Additionally, we introduce a comprehensive framework that unifies various approaches and outline several promising avenues for improving reinforcement learning efficacy.

Despite DeepSeekMath's~notable accomplishments in quantitative reasoning tasks, its abilities in geometry and theorem proving remain less robust compared to leading closed models. For example, during preliminary experiments, the model struggled with questions involving triangles and ellipses, suggesting that the pre-training and fine-tuning data may be subject to selection biases. Moreover, due to constraints in model size, DeepSeekMath~exhibits less improvement in few-shot settings compared to GPT-4; while GPT-4’s performance increases with few-shot prompts, DeepSeekMath~shows comparable results for both zero-shot and few-shot tasks. Moving forward, we plan to further enhance our data selection pipeline to develop a more refined pre-training dataset. We also intend to investigate the potential research directions outlined in Section \ref{sec:effec_rl} to achieve more effective reinforcement learning strategies for LLMs.

\end{document}